\problem{}
Consider the following puzzle set in the classrooms of the School of Information Science and Technology (SIST) at ShanghaiTech University. The grid represents the classrooms, where each classroom may be either empty, occupied by a Computer Science (CS) student, or occupied by an Electrical Engineering (EE) student. The objective is to remove certain students from the classrooms in such a way that the remaining students satisfy the following conditions:

\begin{itemize}
    \item Each row must contain at least one student.
    \item No column can contain both CS and EE students.
\end{itemize}

In some cases, for a given initial arrangement of students, it is impossible to achieve this goal.

\textbf{Prove that determining whether this puzzle is solvable is NP-complete.}

\noindent \textbf{Proof:}

We want to decide whether it is possible to delete some students so that
the final configuration satisfies the above constraints. 
Call this decision problem \(\mathsf{PUZZLE} \).

\noindent\textbf{Step 1: \(\mathsf{PUZZLE} \in \mathrm{NP}\).} \\

A certificate consists of the set of students to be deleted.
A verifier checks in polynomial time that:
(i) after deletions, each row contains at least one student;
(ii) no column contains both a CS and an EE student.
Thus the problem is in NP.

\noindent\textbf{Step 2: Reduction from 3\textsc{SAT} to \(\mathsf{PUZZLE}\).} \\

We reduce from the NP-complete problem 3\textsc{SAT}.  
Given a formula
\[
\varphi = C_1 \land C_2 \land \cdots \land C_m ,
\]
with variables \(x_1,\dots,x_n\), each clause containing exactly three literals,
we construct an instance of \(\mathsf{PUZZLE}\) as follows.

\paragraph{Variable rows.}
For each variable \(x_i\), create one entire row in the grid.
Place a CS student in every cell of this row that corresponds to a literal \(x_i\),
and an EE student in every cell corresponding to a literal \(\neg x_i\).
These two types appear in disjoint columns described below.

\textbf{Clause columns.}
For each clause \(C_j = (\ell_{j,1} \lor \ell_{j,2} \lor \ell_{j,3})\),
create a column in the grid.
In the row of variable \(x_i\), the cell in the column of \(C_j\) contains:
\begin{itemize}
    \item a CS student, if \(\ell_{j,t} = x_i\) appears in the clause;
    \item an EE student, if \(\ell_{j,t} = \neg x_i\) appears in the clause;
    \item empty, if \(x_i\) does not appear in \(C_j\).
\end{itemize}

\paragraph{Intuition.}
A column is conflict-free (no CS/EE mixture) if and only if
we delete all rows (variables) whose literal in that column
disagrees with the chosen truth assignment of the variable.
Since each row must retain at least one student somewhere,
this forces us to pick a consistent assignment for each row.

\paragraph{Polynomial-time.}
The grid has \(n\) rows and \(m\) columns. Filling each cell is constant work.
Thus the construction runs in polynomial time.


\noindent\textbf{Step 3: Correctness (equivalence).} \\

\noindent
\(\Rightarrow\)
Suppose \(\varphi\) is satisfiable, with satisfying assignment \(x_1^*,\dots,x_n^*\).
For each variable row \(i\):
\begin{itemize}
    \item If \(x_i^* = \mathrm{True}\), delete all EE students in that row.
    \item If \(x_i^* = \mathrm{False}\), delete all CS students in that row.
\end{itemize}
Each row keeps at least one student because each variable appears in some clause.
Now consider any column corresponding to a clause \(C_j\).
Since the clause is satisfied, at least one literal \(\ell_{j,t}\) is true.
The student corresponding to that literal remains in the column,
and all students corresponding to false literals are deleted.
Hence the column contains only one type (CS or EE), so no conflicts remain.
Thus this is a valid solution of \(\mathsf{PUZZLE}\).

\smallskip
\noindent
\(\Leftarrow\)
Now suppose the \(\mathsf{PUZZLE}\) instance has a valid deletion pattern.
Since each row retains at least one student,
the remaining students in that row must all be of the same type
(otherwise some column would contain both CS and EE).
Thus each row uniquely determines a truth assignment for its variable:
CS means the variable is assigned True, EE means False.
Let this induced assignment be \(x_1^*,\dots,x_n^*\).

Now consider any clause column \(C_j\).
The column is conflict-free, so all remaining students in it have the same type.
Since the column had at least one student remaining
(otherwise the corresponding clause would have all its rows empty,
contradicting the row constraint), the surviving student corresponds
to one literal of \(C_j\) that is consistent with the assignment.
Thus that literal is true under the assignment.
Therefore each clause has at least one true literal.

Hence the induced assignment satisfies the formula \(\varphi\).

Therefore,  \(\mathsf{PUZZLE}\) is NP-complete.
\newpage